\documentclass[10pt]{article}

\usepackage[T1]{fontenc}
\usepackage[utf8]{inputenc}
\usepackage{changepage}
\usepackage{amsmath}
\usepackage{amsfonts}
\usepackage{amssymb}
\usepackage{url}
\usepackage{indentfirst}
\usepackage{hanging}
\usepackage{titlesec}
\usepackage{hyperref}
\usepackage[notes, short, backend=biber]{biblatex-chicago}
\usepackage{csquotes}

\addbibresource{./right-wing-conversion-narratives.bib}

\setlength{\parindent}{0pt}
\setlength{\parskip}{0.5\baselineskip}

\renewcommand{\labelitemi}{\(\hookrightarrow\)}
\setlength{\itemsep}{0pt}
\titleformat{\section}{\bfseries\Large}{\thesection}{1em}{}
\titleformat{\subsection}{\bfseries\large}{\thesubsection}{1em}{}

\title{right-wing conversion narratives}
\date{2026-09-26}
\begin{document}
\maketitle

Okay, so I was reading Kathleen Blee's \citetitle{Bl02}, which is a sociological book about racist political activists in the United States devoted to investigating the question posed in its first chapter: \textit{what are racist activists like?} In her research for this book, Blee identified and interviewed thirty-four women members of various racist political activist groups (see here\autocite[198]{Bl02} for the methodology) and here presents her interpretations. One of the themes common within interviews was the presence of a conversion story, that is, an answer to the questions, \textit{how and why does one become a racist activist?}

Conversion story is a literary genre that, to ignore most of its long history, is usually associated in the US with the Protestant evangelical movement, and more broadly comprises writing, usually quasi-autobiographical, presenting a narrative account of one's religious conversion. This narrative structure is often applied to stories where one converts not to a particular religion but to some other, more abstract and variably legitimate social movements (like psychoanalysis, homeopathy, veganism, organized misogyny, white nationalism, transphobic activism, or any number of conspiracy theories). The interplay between right-wing groups and the evangelical movement in the United States that characterized the twentieth century Republican Party especially makes the conversion narrative a natural feature of racist autobiography; but it also calls one's attention to the details of that comparison. Religious conversion stories are rarely popular literature outside of evangelical circles and I'm honestly mostly familiar with them by the analogous genre within right-wing politics, so I want to see if this comparison might be made more detailed. 

Blee classifies the groups of which her interviewees are members into four categories: neo-Nazis, racist skinheads, the Ku Klux Klan, or the Christian Identity movement. I'll discuss the distinctions between these groups when necessary, but given the diminishment of the skinhead culture in the twenty-first century, the emergence of myriad other hate groups, and the centrality of the hate movement within the Republican Party, I will more often fall back on the generality of ``the right-wing.'' Although colloquially such politics might be described as ``conservative,'' following Sara Diamond\autocite[6]{Di95} I use the terms ``right-wing'' or ``rightist'' instead.

I'm also not really thinking about narratives of people like David Horowitz or Irving Kristol, rightist public figures who emphasize their prior identification as socialists (sometimes quite minor and sometimes very brief) in their personal branding as political thinkers. While these share some structural aspects with the conversion narratives to follow, they function more to suggest their author's familiarity with the Left or to establish a pretense of being ``politically moderate.'' The purported former leftist narrative remains an irritating and influential trope of political writing in the United States, and I would direct the interested reader to George Packer's 2016 article \citetitle{NewYorker16}.\autocite{NewYorker16} 

\section*{{\itshape a stranger to jesus, a rebel, a thief, a drunkard, a hard drug taker, an adultress, a hippie, a self-centered, confused young woman}\autocite[43]{Br91}}

The conversion narrative genre has a complex history, with tropes and features continually coming into or going out of fashion in response to religious trends. Blee's main reference for this is Virginia Brereton's \citetitle{Br91}, which traces the history of conversion life writing from early nineteenth-century American Protestants, themselves principally inspired by seventeenth- and eighteenth-century English Puritan literature, particularly didactic texts such as Philip Doddridge's 1745 \textit{On the Rise and Progress of Religion in the Soul}.\autocite[10]{Br91} Quickly, here's my impression.

These nineteenth-century narratives tended to present their narrators' lives chronologically, in a predictable five-act structure.\autocite[6]{Br91} The first act usually depicts the narrator as a child and, more abstractly, as unconcerned with or ignorant of the question of salvation. Here, the narrator offers some background biographical information, often setting themselves apart from a more devoutly religious family (of which they are almost always a member), before proceeding through their early life, which is marked by behaviors judged retrospectively by the narrator as sinful or debauched (for women, these could have been ``novel reading, dancing, and theater going'' or their ``quick temper, willfulness, flightiness, or impetuousness,'' while for men, these were more often things like blasphemy and drunkenness). Act Two requires the narrator to become personally concerned with their own sinfulness and salvation. Sometimes precipitated by death, illness, or some other tragedy and sometimes incited by a new acquaintance or reproach from an authority figure, this concern might lead to one or more reform attempts, which necessarily fail due to the narrator's lack of piety. 

Such piety would only emerge in Act Three, which Brereton describes as the ``conviction of sin,'' in which narrators ``would realize that their very wills were corrupt, that they carried internally a \textit{disposition} to sin about which they could do nothing.''\autocite[7]{Br91} Were conversion to take place in a single moment, as many narrators described it, it would be that between the third and fourth acts -- such an instant could coincide with a chance or commonplace encounter between the narrator and a scriptural passage, a friend, or a stranger, with the narrator's response demonstrating the completion of their conversion. In Act Four the narrator would assert a permanent change in their character and would reorient their lives towards evangelism, spreading word of their conversion within the conversion narrative itself as well as joining missionary societies or other formalized religious organizations. The fifth act would comprise the narrator's minor doubts about the movement to which they have converted or about their conversion specifically, but would emphasize the transience of these doubts, and would sometimes end with a posthumous epilogue in which a witness would testify as to the narrator's tranquility and devotion in the moment of their death.\autocite[9]{Br91} 

Throughout the twentieth century, this stock narrative would change in both structure and style. Some modernizations reflected changing social norms -- childhood impetuousness was no longer as shocking a sin, and discussion of sexual topics became more permissible -- and others reflected changes in religious practices -- the concept of original sin was declining in popularity, while converts became more likely to describe ``accepting'' Christ than the prior formula of ``being accepted by'' him.\autocite[56]{Br91} Depicting the instant of conversion as a moment of revelation, already optional in the nineteenth century, would become even rarer, as would the pre-conversion period of ``conviction'' (although the word itself persisted, used in the sense that one would use ``accusation'').\autocite[53]{Br91} On the other hand, the narrator's participation in religious organizations in the fourth act would become more central, probably due to the increasing number and popularity of lay organizations across the country.\autocite[57]{Br91} 

Alongside larger social shifts, evangelicals were faced with the problem of saturation. Nineteenth-century conversion stories were often highly formulaic, relying on stock phrases and scriptural allusions, lending them an impersonality that encouraged nonspecific readers to transpose themselves onto their protagonists.\autocite[18]{Br91} In the twentieth century, though, this style failed to enjoy popularity both among evangelicals (who found them repetitive and unoriginal) and among outsiders (who found such language unfamiliar) -- a problem evangelicals perceived as a hindrance to the movement's ability to recruit new converts, especially as religious conversion itself became less normative.\autocite[73]{Br91} In response, authors adopted stylistic affects from various other genres, writing about religious conversion in legalistic or scholarly terms or through the language of business, sports, or advertising.\autocite[75]{Br91} 

\section*{like if joseph campbell was into racism in a different sort of way}

Contemporaneous with this diversification of style was the evangelical movement's consolidation as a quasi-grassroots political force. ``From the 1940s through the 1960s,'' Diamond writes, evangelicals engaged in prepolitical activity that followed two general patterns. First, they formed parachurch ministries and lobbied against government obstacles to their propagation of the gospel. Second, at the ideological level, they played a role in the Cold War, through publicly accusing liberal church representatives of being communist sympathizers.''\autocite[95]{Di95} The first pattern would enable relatively independent evangelical churches to coordinate nationally via religious magazines, talk radio and, later, televangelism. The second pattern probably arose from some combination of happenstance and natural affinity, and led evangelicals (as opposed to many mainline Protestants at the time) to find anticommunists to be political allies. Founded in 1956, Billy Graham's magazine \textit{Christianity Today} would publish a series of tirades by J. Edgar Hoover, featuring such hits as:
\begin{displayquote}
    The Communists are today spraying the world with ideological and propaganda missiles designed to create a deadly radioactive cloud of Marxism-Leninism.

    The deadliest of these Communist missiles -- whose warheads are exceptionally heavy -- are being directed against the Christian pulpit. Communist gunners, with special ideological training and schooled in atheistic perversity, are ``sighting in'' the clergy -- hoping to shatter, immobilize, and confuse this powerful form of idealism, morality, and civic virtue.\autocite[101]{Di95}
\end{displayquote}

The calcification of the Christian Right as a discrete political force in American partisan politics would, to oversimplify Diamond's analysis, begin somewhere around Barry Goldwater's 1964 presidential campaign and be completed somewhere around the 1980 election of Ronald Reagan.\autocite[172]{Di95} Simultaneously, racist factions of the right wing had declined in electoral influence after George Wallace's failed 1968 campaign, but persisted in the forms of scattered Ku Klux Klan chapters, various tax protestor and militia groups, and a few centralized activist organizations like Willis Carto's neo-Nazi Liberty Lobby.\autocite[257]{Di95} After 1980, these groups would see a resurgence in membership and influence that would lead them back into electoral politics and, importantly for us, into a coalition with right-wing evangelicals. And it's here, among the spawn of this unholy union, that we find the right-wing conversion narrative.

\subsection*{act I: {\itshape Although she acknowledged that she did not see the driver who hit her, Judy nonetheless maintained that it ``must have been'' an African American man.}\autocite[41]{Bl02}}

The first act remains similar. The narrator recalls their wretched life pre-conversion, usually in the sense of something commonplace of which they claim to be a victim. For the less subtle racists, victimization at the hands of racial minorities is immediately and directly the cause of the narrator's sudden turn to a hate group, but others will linger on this point, portraying themselves as tolerant and blameless. Frequently the story begins in childhood, recontextualizing memories of mistreatment within the narrator's current racist worldview while also arousing pity in the audience. 

Helpfully, Blee doesn't dwell on the factuality of her subject's narratives -- the racist activists she interviews are so blatantly untrustworthy that it is more academically useful to analyze their stories as autobiographical fiction. This isn't to say she believes them:
\begin{displayquote}
    Racist activists tends to be disingenuous, secretive, deliberately intimidating, and prone to evasive or dishonest answers. Standard interviews are often unproductive, yielding little more than organizational slogans repeated as personal beliefs.\autocite[201]{Bl02}
\end{displayquote}
On the contrary, she presents the words of her interviewees as works of life history, inherently factitious in the way that any autobiography is. She's also conscious of the academic risk undertaken in forming relationships with such people, and of the care that must be taken to prevent her subjects' political propaganda from communicating itself via her work. Her approach to this problem is interesting and emotionally complicated and I really do recommend reading her introduction because of it, but she tries in spirit to assume most of the details in her interviewees' stories are fabricated:
\begin{displayquote}
    I use incidents and information from the life histories as data only when these are contrary to the overall impression that a woman is trying to convey (e.g. admissions of accepting welfare support) or are supported by additional details from the structured questionnaire and chronological life account.\autocite[203]{Bl02}
\end{displayquote}
While this method might filter out the most egregious falsehoods in her subjects' narratives, it must rely on her own judgments of those subjects' (usually nefarious) intentions. On the other hand, it precludes her from really spending that much time on the way those falsehoods function in service of those intentions. 

Anyway, Blee's subjects tend to interpolate historical references to the racial integration of schools, which would have been at least chronologically accurate to their lives, with a discourse of school bullying contemporary to the 1990s. Such historical reference points exist not really to make these anecdotes more believable but to bridge a divide between the narrative expectations of the racist fringe and those of the (not un-racist) white mainstream. We can imagine a more credible-sounding (while still obviously warped) story, one where the narrator as a kind-hearted child finds themselves befriended by, then betrayed by other children -- but such a story would demand an attention to non-white characters that racists would find unacceptable. Narrators instead construct parables where the vagaries of their conflict with an amorphous racial other are masked behind widely-known sociohistorical themes, where the audience's familiarity can be expected to compensate for such poorly-developed stories. 

% absolutely made up yet commonplace grievances

% also childhood memories

\subsection*{act II: {\itshape ``Someone gave me a tip one time as they were leaving. And inside the five dollar bill was a card that said, `You have been patronized by the KKK.' You can't tell. It's real surprising.''}\autocite[4]{Bl02}}

Encounters with recruiters for racist groups are played as coincidental and commonplace. The chance encounters of evangelical conversion narratives, insinuating some small yet divine intervention in that narrator's life, here suggest an omnipresence of rightist activism. Here, Blee presents two perspectives. In her subjects' stories, ``simple happenstance is often an element of racist affiliation--'' that is, many describe their paths to joining a racist organization as beginning with a brief acquaintance, reconnecting with an old friend, or coming across a piece of minor propaganda.\autocite[30]{Bl02}

Such an illusion of spontaneity is undercut by the well-preserved manuals and testimony of those recruiters, which detail a strategy of identifying vulnerable targets and slowly befriending them before persuading them to join a hate group. In the words of German neo-Nazi Ingo Hasselbach,
\begin{displayquote}
    I liked to approach 14- to 16-year-olds after school. Usually they didn't really have a political position, but for whatever reason they'd decided it was cool to be right-wing. The first thing I did when I met one of these boys was to show that I wanted to be his friend. I'd always slip in a bit of ideology against foreigners along the way, but only casually at first.\autocite[28]{Bl02}
\end{displayquote}
Combining these perspectives, it's tempting to see recruitment into racist groups as a process of indoctrination, in which the recruiter, unbeknownst to the recruit, manipulates them into joining. This would, I think, be putting too much faith in Blee's interviewees' honesty. 

Rather than these recruitment efforts appearing to the unsuspecting as random encounters, it's more productive to read them as self-effacement. In the corresponding parts of twentieth-century evangelical narratives, ``the sinful \dots\ often betray a trouble-seeking, aimless restlessness,'' enticing the identification of a hypothetically aimless reader -- the ideal convert.\autocite[50]{Br91} The rightist narrative too desires this identification, but follows it with a sleight of hand. For religious converts, it is essential to locate the moment of conversion within a personal choice, but for rightists this would be counterproductive. More fundamentally, as a political argument, proclaiming one's decision to be racist really just invites confusion and suspicion and would undermine the implied \textit{inevitability} of the conversion.

%\subsection*{act III: {\itshape ``I wish we could get what we want through negotiation, but I don't see that happening. I see it all coming down to civil war.''}\autocite[181]{Bl02}}

\subsection*{act III: {\itshape ``I would say a race war is inevitable, but I couldn't say when.''}\autocite[98]{Bl02}}

Such a claim to inevitability is implicit in most right wing beliefs, particularly those derived from the immutable hierarchies rightists imagine shape human society. Like, thinking of feminism or civil rights as wastes of time because they challenge social structures they believe fundamentally unchangeable, that sort of thing. An anxiety about damnation that became in the twentieth century an anxiety about a lifetime unfulfilled realizes itself here as the millenarian fantasy of a race war. For white supremacists, imagining racial equality to be not just undesirable but impossible causes them to imagine a future race war a certainty, thereby dodging the problem of choosing to be white supremacists. If a race war is inevitable, they say, why not become a racist activist? Otherwise, aren't you just wasting time?

The exact flavor of millenarianism varies, as it does among evangelicals: some ``see it all coming down to civil war;''\autocite[181]{Bl02} for others, ``the race war is here and has been going on since biblical times.''\autocite[99]{Bl02} In either case, the narrator frames their participation as unavoidable and, in doing so, implies that \textit{anyone's} participation is unavoidable. Internally, this belief in some immanent war intermittently motivates acts of violence against the group's perceived enemies, but it persistently motivates and is motivated by what Blee describes as ``narrative violence,'' that is, ``violence as tradition and commonsense understanding, something taken for granted, seen as natural.''\autocite[177]{Bl02} In this segment of the conversion story, the narrator describes the process by which they became accustomed to violence through the group's own violent mythology.

The narrative violence described here can take several forms. Blee's interviewees from skinhead groups often recount hazing rituals or fights commonplace in the white power music subculture; one Klanswoman recalled being pelted with rocks while participating in a racist parade; others describe group discourses about violence and its justifications. Critically, the explication of narrative violence serves to accustom the audience alongside the narrator to that violence. Possibly due to there being relatively fewer acts of strategic violence, or (in my opinion, more likely) due to a perceived greater transgressiveness of strategic violence, narrators deliver their life stories against a backdrop of unremarkable violence so that the racist violence their group perpetrates might too seem unremarkable. It is for this reason that Blee says, ``narrative violence must be seen as a precursor to strategic violence.''

``Like the raw celebration of death and violence that,'' she continues, ``characterizes fascist movements and organized crime, the violent culture of organized racism has a political effect, however remote from any particular action.''\autocite[186]{Bl02} Although they usually lack the means to commit the world-shaping acts of violence about which they fantasize, racist movements' narrative violence shouldn't be seen as a substitute or impotent pretense. It's more a kind of aspirational necropolitics, an act by which racist activists performatively abject themselves from the state and its legalism. Achille Mbembe describes the colony as ``the location par excellence where the controls and guarantees of judicial order can be suspended -- the zone where the violence of the state of exception is deemed to operate in the service of `civilization.'''\autocite[24]{Mb03} This is essentially the dream of fringe rightist activism: personally lacking the sovereign political power of the state and frustrated by its perceived unwillingness to exercise its right to kill, activists seek to construct a space -- a colony -- where that right is exercised for its own sake. In this dream, the inevitability and necessity of the race war can persist, no matter how bizarre or implausible; ``here, the fiction of a distinction between `the ends of war' and `the means of war' collapses.''\autocite[25]{Mb03}

% racist millennialism. unlike in the nineteenth century, these are universally hopeful /goals/ rather than a fear that incited the conversion. 

% we are missing something about anxiety, and the invented patheticness

\subsection*{act IV: {\itshape ``The `new world order' is the worst thing ever, especially for the white race and also for every single individual in the world --- besides the government and Jews.''}\autocite[94]{Bl02}}

If in Act III the convert synchronized their system of beliefs with the eschaton they have come to expect, in Act IV they apply these principles to their own life. Narrators tend to signal the completion of their conversion in two ways: reinterpreting their personal grievances into established conspiracy theories and formally joining a rightist organization. As right-wing groups typically both require belief in and produce idiosyncratic conspiracy theories, these two acts are inseparable. 

In one direction, a recruit to a racist group has to familiarize themselves with labyrinthine and esoteric racial beliefs. Blee describes this familiarization as a transformation of ``everyday racism'' into ``extraordinary racism'':
\begin{displayquote}
    Everyday racism, which reflects negative, stereotyped views of minority groups, lacks any systematic explanatory force. It is pernicious but unremarkable among whites in a racist society. In racist groups, however, \dots\ it is molded into extraordinary racism, an ideology that interprets and gives meaning to a wide variety of phenomena that seem unconnected to race, ranging from the global economy and the growth of media monopolies to more immediate personal issues such as the quality of family life, city services, and medical care. \dots\ The difference between everyday and extraordinary racism is the difference between being prejudiced against Jews and believing there is a Jewish conspiracy that determines the fate of individual Aryans, or between thinking that African Americans are inferior to whites and seeing African Americans as an imminent threat to the white race.\autocite[75]{Bl02}
\end{displayquote}

One could assume, then, that \textit{anyone} with everyday racist beliefs could, with some study, become an extraordinary racist. This, while true in a way, is in another way an assumption exploited in racist proselytization. It suggests that there is a course of study, starting from widely-held beliefs, which invariably leads to extraordinary racism. However, this would fail to grasp conspiratorial belief as rhetorical affect, to assume incorrectly that racist activists are persuaded into their extraordinary racism. In fact, racist propaganda delivered directly is not especially persuasive, even to racists:
\begin{displayquote}
    At a neo-Nazi gathering I attended, most people paid only sporadic attention to long, boring speeches by the group's self-proclaimed leaders. Even a livelier (at least to me) presentation by two younger members \dots\ had no more success in sustaining the interest of the audience, many of whom left early or spent time conspicuously reading the newspaper.\autocite[76]{Bl02}
\end{displayquote}
Racist literature and broadcasting don't fare much better:
\begin{displayquote}
    The rhetoric of almost all written racist propaganda certainly presents a formidable obstacle to any but the most dedicated reader. Books, newsletters, newspapers, and magazines published by most racist groups are written in an impenetrable style. Long blocks of text justifying the group's positions are broken occasionally by crude cartoons or caricatures of enemies, mixed with references to obscure internecine battles among racist leaders.

    Almost all the women I talked with were aware of racist TV shows, carried from time to time on a hodgepodge of local cable access channels. Yet \dots\ only two were able to name the featured guests or the topic of more than a single broadcast. A Klanswoman told me that one of them is ``the damn silliest thing ever and [its host] is a total buffoon.''\autocite[77]{Bl02}
\end{displayquote}

Even if competently presented, the beliefs of extraordinary racism tend to come across as ridiculous, as flimsy networks of verifiable falsehoods connected by illogical imaginations. Despite its unconvincingness as argumentation, though, extraordinary racism functions as a rhetorical style, with familiarity being expected of and necessary for members of rightist organizations. Learning these aesthetics of extraordinary racism serve specific purposes both for the practices of conversion and for narratives of conversion. 

For conversion as practice, the language of extraordinary racism transcends the boundaries of individual organizations, making fluency essential for intergroup communication and recruitment. Allusions to shadowy forces form a stylistic bridge between mainstream politician Ron Paul, speaking at a 1994 meeting of the ``paleoconservative'' Randolph Club on ``some of the `real' victims of state authority, [including] anyone who won't pay their taxes or obey `handicapped only' parking signs,''\autocite[182]{Di96}, and white supremacist Jared Taylor, delivering at the same conference what Sara Diamond called ``one of the wildest racist diatribes I can ever recall hearing'' to the same, totally blas\'e audience.\autocite[183]{Di96} Similarly, the 1995 ``Freedom Rally,'' a conference of about 200 right-wing activists from across the country, featured everything:
\begin{itemize}
    \item The National Commodity and Barter Association, a precious metals-based, Posse Comitatus-affiliated tax evasion scheme;
    \item Militia organizer Eugene Schroder, alleging how a ``secret government dictatorship'' was set up during the Roosevelt administration;
    \item Christian Identity pastor Everett Sileven, advocating for uncertified religious schools, hawking a pamphlet called ``The Sin for Which God Will Kill: Inter-Racial Fornication by Copulation and Reproduction,'' and giving a speech denouncing the ``wickedness [of] the American monetary system, the income tax, race-mixing and multiculturalism, government schools, sodomy, women's lib, and the United Nations;''
    \item sovereign citizen Godfrey Lehman (who Diamond notes partially because he is Jewish), advertising a new militia group in Berkeley, California; 
    \item Bay Area Klansman and local write-in candidate Richard Charles; and
    \item extensive merchandise for sale, including militia training manuals, some sort of homebrew caffeine supplement, instructions on how to file fraudulent tax documents, new editions of the notorious racist tabloid \textit{Jubilee}, books like \textit{The Biological Jew} by white supremacist Holocaust denier Eustace Mullins, and some classics, like audio tape translations of Hitler's speeches.\autocite[210--214]{Di96}
\end{itemize}
The particular organizations and individuals represented at the ``Freedom Rally,'' while espousing similarly racist beliefs, aren't really interchangeable and frequently differentiate themselves in their tactics, membership, and political priorities. Even just the groups that Blee investigates (the KKK, neo-Nazis, skinheads, and the CI movement) have substantively different practices from each other, but they are all enabled to participate in cross-organization activity and discourse by the mutually-intelligible language of extraordinary racism. The conversion narrative is a site where such a language can be practiced, reproduced, and -- critically -- taught.

For conversion as narrative, the same language operates in part as specialized jargon does for any secret knowledge grift. The aesthetics of extraordinary racism suggest mystery and evoke curiosity while simultaneously miring outside observers in nonsensical, easily changeable mythology. Counterintuitively, it can often be to the political benefit of organized racists to cite as motivation for their activism convoluted bullshit, like the perennial ``cultural Marxism,'' ``white genocide,'' or ``Zionist Occupation Government'' conspiracy theories, because to do so tempts outsiders into inferring the racists' goals to be equally detached from reality. In fact, racist activists have simple and comprehensible political goals that they frequently achieve: to disenfranchise, subjugate, or otherwise harm racial minorities, in whatever circumstance, given whichever opportunity.

\subsection*{act V: {\itshape ``The picture I had in my mind was a lot different than it finally turned out,'' commented an eighteen-year-old who had moved to Idaho to enter into a polygamous relationship with a white supremacist.}\autocite[153]{Bl02}}

Potentially unique to the rightist conversion narrative is a belief that conversion has made the narrator's life strictly worse. The salvation accorded evangelical converts is definitionally an improvement in their circumstances, and the narratives of other movements Brereton considers -- Alcoholics Anonymous, feminism, and, unfortunately because it was the nineties, Weight Watchers -- all share the expectation that the narrator post-conversion is both psychically and materially better off than they were pre-conversion. Contrastingly, right-wing converts tend to emphasize a burden that comes with their newfound awareness: this is a staple of all secret knowledge grifts, with the intent to inflate the magnitude of the unconvincing, nonsensical, or widely-available knowledge they claim to hold. Alongside this burden is rightists' consciousness that their views are frequently so antisocial, bigoted, or otherwise offensive as to essentially always come with some degree of social ostracism. Among Blee's sources, such a consciousness is peculiar to racist activist women, who tend to feel that ``it's painful, it hurts, it's all-consuming when you have the knowledge'' -- the ``knowledge,'' of course, being white supremacist beliefs.\autocite[47]{Bl02}

Separate from this absolutely laughable whining about the psychic burden of being a white supremacist, and from the social consequences brought upon those ``dumb enough to walk down the street with a swastika on their T-shirt'' (in one's words), are the gendered consequences racist women predictably suffer. Right-wing politics tends to offer women very little, with its fetish for authoritarian, male-dominated hierarchies -- an ostensible contradiction that has been studied extensively. More surprising than racist women's participation in organizations explicitly opposed to women's rights are the conflicts over sexism that take place within those organizations. Although right-wing organizations are generally explicitly antifeminist, and this antifeminism easily lends itself to racist and antisemitic conspiracy theories (``feminism is the means to weaken Aryan masculinity, promoted by the international Jew,'' according to one 1990s pamphlet and innumerable 2020s podcasters\autocite[145]{Bl02}), it is relatively common for women and some men to discuss the problem of gendered violence within their groups. It is not lost on them that the violent abuse of women is simultaneously an absolute hindrance to their movement and the only conceivable outcome of a violent movement predicated in part on abusing women. 

Criticism of the movement can then become a reaffirmation of the movement's ideals. Members juxtapose expressions of faith in the goals of the movement (establishing white separatist communities, mobilizing other white people into hate organizations, etc.) with a competitive distrust of the \textit{individuals} in the movement, constantly criticizing each other for their perceived lack of commitment to the cause, and often -- correctly -- identifying each other as, for example, ``scumbags with fraudulent, unmotivated behavior.''\autocite[50]{Bl02} For a multitude of probably obvious reasons, rightist groups are, like evangelical ones, rife with scams -- members have by virtue of their membership already announced themselves as unskeptical, prone to fantasy, and willing to stake their finances and reputation on anything that interpellates their particular beliefs. 

Actual attempts to address some (but only ever some) forms of misogyny within right-wing groups do exist to various extents, but conversion narratives that engage with them are usually not optimistic. Given the hostility towards dissent within rightist groups, the conversion narrative provides an approved space for this discourse to occur -- in much the same way that historical conversion narratives constituted a rare (at least within American Christianity) medium where women's public participation could be made acceptable.\autocite[93]{Br91} The external purpose of such criticism is more complicated. In one sense, it could be meant to soothe a listener's misgivings or performatively engage in insincere reflection. In another, it discourages potential converts unlikely to tolerate the personal harm group members are almost certain to suffer. In yet another, it appropriates the language of mainstream feminism in a gleeful mockery, disguising rightist politics within a socially acceptable discourse and using that disguise to insinuate the group's legitimacy. 

% okay i cut this section but i'm leaving it here bc if you find this you HAVE to read g gordon liddy's book.
%
%\section*{the tv movie adaptation: {\itshape\large ``Viewers may justifiably wonder what possible purpose has been served.''}\autocite{NYT82}}
% one extreme is more political, G. Gordon Liddy. one extreme is less political, incels. so far we have been centrists; thetheir y are political racists but not politically active like g gordon liddy was
%
%In 1980, Watergate burglar G. Gordon Liddy, his twenty-year prison sentence just commuted by Jimmy Carter, slapped together \textit{Will}, an autobiography calculated to launch his career in right-wing talk radio. That's viewers of \textit{Will}'s TV movie adaptation in the section title, according to the New York Times. I don't wonder about the purpose that book served, even though I do sometimes wonder whether it could ever possibly be worth it, reading so much neo-Nazi propaganda. 
%
%Contemporaneous reviewers at least understand that Liddy ``seems to exist in a world of absurd macho fantasy,'' and that ``his basic social and political instincts are fascistic.''\autocite{NYT82} In 2026, old men at New York bars remember more the absurd macho fantasy of it, like his desperate anxiety to remind the reader that he owns a gun or his widely-parodied candle self harm bit, but I cannot possibly overstate how much this book is about G. Gordon Liddy loving Hitler. 
%
\section*{recidive: {\itshape ``Therapy is a colossal waste of time. What, women will suddenly find me attractive \dots\ because I go to therapy?''}\autocite[70]{Ke24}}
 
When I said the phrase ``right-wing conversion narratives'' to my roommate the other day, he immediately was like, ``oh, like incels? `I took the red pill,' that sort of thing?'' The recognizability and proliferation of incel narratives, probably disproportionate to the number of adherents, have finally caught the attention of academics, but just barely. In the common understanding, ``incel culture'' is a right-wing reactionary movement that aims to punish, humiliate, and subjugate women. Analyses often focus on its origins on the internet, its strange fixation on seemingly trivial bodily insecurities, its codification of truly absurd and usually racist hierarchies of physical attractiveness, and its propensity to foment and precipitate mass violence.\autocite[68]{Ke24} There's an understanding both quantifiable\autocite[in][the authors compare a body of canonical white supremacist texts with masculinist ones, using statistical techniques to detect overlapping topics and sentiments.]{PLPV22} and anecdotal\footnote{my roommate, personal communication with the author, 2026.} that incels as a group overlap extensively with white supremacist, neo-Nazi, and other right-wing hate groups. 

{%
\vspace*{\parskip}
\vbox{
    Apparently, some disagree. A frequently-cited study by Costello et al. of 151 self-identified incel men\stepcounter{footnote}\footnotemark[\thefootnote] asserted no statistically significant differences in \small political orientation when compared to non-incels, leading some reporters to make the claim, surprising given stereotypes, that 44.7\% of the surveyed incels were left-leaning compared to only 38.85\% right-leaning. Such reporting is regrettable \footnotesize -- first, the study determined political orientation by a single test item, in which participants were asked to place themselves on a five-point scale, with 1 being `left wing'' and 5 being ``right wing,'' so we can at best take their responses to be indicative of their own self-professed, possibly poorly understood political identities. \scriptsize Second, participants for the study were recruited from -- and with assistance from the administrators of -- incel internet forums, so we should view their responses as narrative building, as a presentation of themselves \tiny curated for outsiders to observe. Third, and most importantly, who cares where on the five-point political spectrum incels place themselves? They proliferate manifestos on the internet ad nauseam -- that's their whole thing! -- and one should be capable of interpreting the political implications of such.
\stepcounter{footnote}\footnotemark[\thefootnote]%
}%
\addtocounter{footnote}{-1}
\footnotetext{\cite{CRTS22}}%
\stepcounter{footnote}
\footnotetext{%
I still should explain why incel beliefs, as they present them, are categorically right wing. So let's clarify what, I guess, if we have to, what right wing movements are. In \textit{Roads to Dominion}, Diamond identifies three themes, around some combination of most explicitly political rightist movements in the US organize themselves: so-called ``free market'' capitalism; anticommunism or, more generally, US military hegemony over the rest of the world; and ``traditional morality'' or, really, white male supremacy. Although rightist political activities are varyingly pro- and anti-state, these themes are linked by a common belief as to the acceptable role of the state: ``to be right-wing means to support the state in its capacity as an \textit{enforcer} of order and to oppose the state as \textit{distributor} of wealth and power downward and more equitably in society'' (9). 

Incels' conception of the state as it exists is not especially complex -- they believe it is controlled by Jews. More precisely, explicitly political discourse on popular, publicly-viewable incel forums will use ``the Jews'' indiscriminately and constantly as a metonym for any state or any apparatus of power at all. Here, the antisemitic canard of secret Jewish control of the state is stripped to its basest form; incels mimic the extraordinary racism of their rightist compatriots without appearing to share their intentionality (one ``off-topic'' forum thread devoted to Israel's ongoing genocide in Palestine comprised, along with a swath of anti-Jewish racial slurs and plenty of other ethnic hatred, a bored attitude of \textit{support} for Israel). The impression is one of racism as affect, which, as it happens, is not so dissimilar from the ill-informed antisemitism of some of Blee's interviewees. As for the state incels desire, they commonly wish for the state to ``more equitably distribute'' women in society, to grant them and every other lonely little man a government-assigned girlfriend in some glorious communist revolution against the tyranny of so-called sexual market value. One should reject that charade entirely; it's repulsive. To read such a desire as for fair distribution of wealth rather than for the enforcement of a gender-based system of slavery would be shameful, an error that uncritically accepts the total commodification of women within their delusional, totally financialized model of sexuality.

Of course, if you're familiar with incel forums to any extent, you'll notice before any of this the sheer volume of racist, homophobic, transphobic, et al. slurs -- a kind of insistence that I imagine would seem to the average 4chan user a little desperate for attention.
}%
}%

And, like. We're on the internet. I don't know that we need to tediously pick apart how public images of incels are constructed intentionally via the language of conversion narrative. They're notorious for their formulaic and broadly suspect stories of lasciviousness and rejection, recontextualized as invitations to hate women categorically. They fixate on physical characteristics that are at once laughably unremarkable and extraordinarily racist. They signal their sincerity by positing truly unhinged conspiracy theories, engaged in constant internal competitions to believe the most, to be the most devoted, the most miserable. They believe they are owed dominion over not just women but the polluted, the degenerate, the unworthy; they yearn for a violent realignment of society that would grant them this dominion. Some of them conclude that such a violent realignment will only come about if they themselves commit those acts of violence. 

Ignoring the political aspect of organized misogyny would necessarily be to ignore the rhetoric of their texts. Interlopers often gravitate towards studying incels' mental health, finding high incidences of depression, anxiety, bullying, self harm, drug abuse, and so on. Deniese Kennedy-Kollar writes that many ``insist that all their issues stem from simply being unable to get sex, which they believe to be an externally caused problem,'' and suggests that their misogyny provides a convenient excuse to avoid interrogating other causes for their problems, especially those for which they might be responsible.\autocite[70]{Ke24} Probably -- but one easily also imagines ways online hate forums predicated on some degree of self-loathing might not be the most mentally healthy environments. This, of course, is a storytelling device, one not entirely consistent with Blee's experiences among organized racists but entirely consistent with Brereton's observation of Christian writing, that ``what unites all the troubled narrators, besides their pain, is the sense of being unloved and unlovable, of isolation, abandonment, and hatred turned inward.''\autocite[51]{Br91}

Tempting as it is to find the causal role played by these purported experiences, it would be to ignore their narrative role. At one level, they play the same role as an inciting incident in racist conversion stories, a vague personal grievance set atop an evocation of large-scale social discourses: here, incels exploit the reification of mental health discourses in mainstream consciousness. Their willingness to identify with broadly understood mental health struggles is only kind of a ``lack of insight,'' and not primarily even an appeal for sympathy -- it is a performance rehearsed to make themselves externally comprehensible and internally homogenous. 

Like organized racist women, incels are tasked with recruiting for an ideology that is on its face repellent and thoroughly unproductive, one that is certain to bring them ostracism, tribulation, and, they might dream, martyrdom. They imbue their narratives with this spectacular misery that filters out the noncommittal while aggrandizing the committed, augmenting stories of their miserable lives with a suggestive \textit{doesn't it sound a little like yours?} But no, it doesn't. They're bigots, frustrated that their crass bigotry is pathetic and contemptible.

\printbibliography
\end{document}

